\section*{Planteamiento del Proyecto}
\justifying
El planteamiento o definición correcta del proyecto es lo primero que se debe lograr para no
desviar el objetivo ni generar cuestionamientos irrelevantes. El alumno debe ser capaz no sólo
de conceptuar el proyecto sino también de verbalizarlo en forma clara, precisa y accesible,
de manera tal que el lector lo comprenda por el documento mismo.

En esta parte se trata de brindar una descripción concreta del proyecto a realizar,
dando una versión de los hechos y fenómenos cuya explicación debe ser interesante y útil,
tanto para el alumno y el lector como para el medio académico y la sociedad. Con tal fin,
partiendo de lo particular y hasta lo general, se explicará el cuestionamiento y la
problemática que dirigirá el proyecto así como las dificultades y dudas que se pretenden
estudiar.

Para tal efecto, el planteamiento del proyecto se puede hacer a partir de una aseveración,
de la cual posteriormente se derivarán una serie de preguntas centrales, que dan origen
a la investigación y que, definidas con claridad y sin ambigüedades, serán de utilidad
para dirigir el proyecto. 

Para ello, se incluirán los hechos, relaciones y explicaciones que fundamenten la
problemática, mencionando aquellos datos que la puedan soportar, ya sea que se
encuentren en otros proyectos o en teorías ya establecidas, por ejemplo.

El planteamiento del proyecto debe dimensionarse apoyándose en cuadros de estadísticas,
figuras, diagramas, etc. Al término de este apartado, el lector deberá estar
plenamente convencido de que por su magnitud, el proyecto realmente requiere
de un estudio que le aporte soluciones.
